\documentclass[10pt,a4paper]{article}

\usepackage[T1]{fontenc}
\usepackage[utf8]{inputenc}
\usepackage{charter}
\usepackage[a4paper,margin=1.9cm]{geometry}
\usepackage{microtype}
\usepackage{xcolor}
\usepackage{enumitem}
\usepackage{titlesec}
\usepackage{tabularx}
\usepackage{booktabs}
\usepackage{amssymb}
\usepackage{needspace}
\newcolumntype{Y}{>{\raggedright\arraybackslash}X}
\usepackage[hidelinks]{hyperref}

\definecolor{accent}{HTML}{1F4E79}
\definecolor{muted}{HTML}{555555}
\definecolor{rule}{HTML}{C9D3DD}

\setlength{\parindent}{0pt}
\setlength{\parskip}{4pt}
\pagestyle{plain}

\titleformat{\section}{\large\bfseries\color{accent}}{\thesection.}{0.5em}{}
  [\vspace{1pt}{\color{rule}\titlerule}]
\titlespacing*{\section}{0pt}{12pt}{5pt}
\titleformat{\subsection}{\normalsize\bfseries}{}{0pt}{}
\titlespacing*{\subsection}{0pt}{8pt}{2pt}

\setlist[itemize]{leftmargin=1.3em,itemsep=1pt,topsep=2pt,parsep=0pt}
\setlist[enumerate]{leftmargin=1.5em,itemsep=1pt,topsep=2pt,parsep=0pt}
\setlist[description]{leftmargin=6.4em,labelwidth=5.8em,labelsep=0.6em,
  itemsep=1pt,topsep=2pt,parsep=0pt,font=\normalfont\bfseries\color{muted}}

\newcommand{\cb}{\(\square\)\hspace{0.45em}}
\newlist{checklist}{itemize}{1}
\setlist[checklist]{label=\(\square\),leftmargin=1.5em,itemsep=1pt,topsep=2pt,parsep=0pt}

\begin{document}

{\LARGE\bfseries\color{accent} ML System Design Roadmap}\\[3pt]
{\large Four weeks of interview preparation for senior applied scientist and
data scientist roles}\\[2pt]
{\color{muted}Raheel \textbullet{} October 2026 \textbullet{} 5 to 6 hours per week}

\vspace{4pt}{\color{rule}\hrule height 0.8pt}\vspace{6pt}

\textbf{Goal.} By the end of week 4, take any open-ended prompt (``design a
system that \ldots'') and talk through a complete, defensible design in 40
minutes, including the serving and monitoring parts that research-focused
candidates most often leave thin.

\textbf{Weekly rhythm.} Four weekday sessions of 45 to 60 minutes for reading
and concepts, plus one weekend session of about 90 minutes: one timed case
spoken aloud (40 minutes), then a written review against the rubric in
Section~\ref{sec:rubric}. Speaking aloud is the part that matters most; reading
alone does not prepare you for the round.

%----------------------------------------------------------------------
\section{The answer framework}

Use the same eight steps for every case. The minutes are a guide for a
45-minute round and leave a short buffer for questions.

\begin{center}
\begin{tabularx}{\textwidth}{@{}r l Y r@{}}
\toprule
 & \textbf{Step} & \textbf{What to cover} & \textbf{Min} \\
\midrule
1 & Clarify & Business goal, users, scale, latency, what action the output
    drives, constraints (regulation, explainability) & 5 \\
2 & Metrics & One offline metric, one online or business metric, and the
    guardrail metric that stops you over-optimising & 4 \\
3 & Data and labels & Sources, how labels arise, label delay and noise,
    class imbalance, leakage risks & 6 \\
4 & Features & Entity, aggregate, velocity, sequence and graph features;
    which are available at decision time & 6 \\
5 & Model & Simple baseline first, then the stronger model and why; training
    cost and interpretability & 6 \\
6 & Validation & Temporal splits, backtesting, calibration, slice analysis,
    threshold choice against capacity & 4 \\
7 & Serving & Real-time versus batch, latency budget, feature freshness,
    fallbacks, how the score becomes a decision & 7 \\
8 & Monitoring & Data and concept drift, alerting, feedback loop, retraining
    trigger, safe rollout & 5 \\
\bottomrule
\end{tabularx}
\end{center}

Steps 2 to 6 are your home ground from fraud modelling. Steps 7 and 8 are the
gap this roadmap is built to close, and they are also where interviewers
separate senior candidates from mid-level ones.

%----------------------------------------------------------------------
\section{Week-by-week plan}

\needspace{9\baselineskip}
\subsection{Week 1: Framework and home ground}
\begin{description}
  \item[Read] Huyen, chapters 1, 2, 4, 5 and 6 (overview, system design,
    training data, feature engineering, model development and offline
    evaluation).
  \item[Learn] The eight-step framework until you can recite it without notes.
    How to state requirements as numbers: transactions per second, latency
    budget, acceptable false-positive rate.
  \item[Practise] Case~1 (real-time payment fraud) twice: once written and
    untimed, once spoken and timed.
  \item[Done when] You have a two-page written design for Case~1 that you
    would be happy to defend, using only non-confidential detail.
\end{description}

\needspace{9\baselineskip}
\subsection{Week 2: Serving and production (the main gap)}
\begin{description}
  \item[Read] Huyen, chapters 7, 8 and 9 (deployment and prediction service,
    distribution shifts and monitoring, continual learning and testing in
    production). Skim chapters 3 and 10 for vocabulary.
  \item[Learn] Batch versus online prediction; feature stores and
    online/offline skew; latency budgets; shadow deployment, canary release
    and A/B testing; retraining triggers; model fallbacks.
  \item[Practise] Case~2 (account takeover) and Case~3 (lending application
    fraud), both spoken and timed.
  \item[Done when] You can spend seven unbroken minutes on serving and
    monitoring for a fraud system without drifting back to modelling.
\end{description}

\needspace{9\baselineskip}
\subsection{Week 3: Cases outside fraud}
\begin{description}
  \item[Read] Aminian and Xu: the introductory chapter, then the video
    recommendation, ad click prediction and harmful content detection
    chapters.
  \item[Learn] Candidate generation followed by ranking; embeddings and
    two-tower models; approximate nearest-neighbour retrieval; cold start;
    position bias; exploration versus exploitation.
  \item[Practise] Case~4 (recommendation), Case~5 (ad click prediction) and
    Case~6 (harmful content detection).
  \item[Done when] You can design a retrieval-and-ranking system from a blank
    page, since interviewers often deliberately leave your domain.
\end{description}

\needspace{9\baselineskip}
\subsection{Week 4: LLM systems and full mocks}
\begin{description}
  \item[Read] One current engineering write-up each on retrieval-augmented
    generation, agent design, and LLM evaluation (company engineering blogs
    are the best source; pick recent ones).
  \item[Learn] The RAG pipeline (chunking, embedding, retrieval, reranking,
    generation); agents with tools; guardrails; evaluation with fixed test
    sets, model-graded checks and human review; cost, latency and caching;
    audit trails for regulated decisions.
  \item[Practise] Case~7 (fraud investigation assistant) and Case~8 (dispute
    resolution assistant), then two full 45-minute mocks with a peer, or
    recorded and replayed.
  \item[Done when] Both mocks score 3 or more on every rubric line.
\end{description}

%----------------------------------------------------------------------
\needspace{16\baselineskip}
\section{Case bank}

\begin{center}
\begin{tabularx}{\textwidth}{@{}r l Y@{}}
\toprule
 & \textbf{Case} & \textbf{What makes it hard} \\
\midrule
1 & Real-time payment fraud detection & Sub-second latency, extreme
    imbalance, delayed and noisy labels, adaptive adversaries \\
2 & Account takeover detection & Session and device signals, few confirmed
    labels, friction cost of step-up authentication \\
3 & Lending application fraud & Synthetic identities, bureau and alternate
    data, no outcome label for rejected applicants \\
4 & Video or product recommendation & Two-stage retrieval and ranking, cold
    start, feedback loops in the training data \\
5 & Ad click prediction & Very high volume, calibration matters, continual
    training on fresh data \\
6 & Harmful content detection & Multi-modal inputs, policy changes, human
    review capacity as a constraint \\
7 & Fraud investigation assistant (LLM) & Grounding in case data,
    hallucination risk, auditability, analyst trust \\
8 & Dispute resolution assistant (LLM) & Tool use, escalation rules,
    evaluating open-ended answers \\
\bottomrule
\end{tabularx}
\end{center}

For Cases 1 to 3, prepare the points that tie back to your own work: sampling
under extreme imbalance, training with noisy labels, and why random splits
overstate fraud model performance.

%----------------------------------------------------------------------
\needspace{18\baselineskip}
\section{Concept checklist}

Tick each item once you can explain it in two minutes with an example.

\begin{minipage}[t]{0.48\textwidth}
\textbf{Data and features}
\begin{checklist}
  \item Label delay and how it biases recent data
  \item Leakage through future information
  \item Point-in-time correct feature joins
  \item Feature store: offline and online paths
  \item Online/offline skew and how to detect it
  \item Streaming versus batch aggregates
\end{checklist}
\textbf{Serving}
\begin{checklist}
  \item Batch versus online prediction
  \item Latency budget across the request path
  \item Model compression and caching
  \item Fallback when a model or feature is down
  \item Score to decision: thresholds and rules
\end{checklist}
\end{minipage}\hfill
\begin{minipage}[t]{0.48\textwidth}
\textbf{Release and monitoring}
\begin{checklist}
  \item Shadow deployment and canary release
  \item A/B testing and its limits in fraud
  \item Data drift versus concept drift
  \item Monitoring without immediate labels
  \item Retraining triggers and cadence
  \item Feedback loops and selection bias
\end{checklist}
\textbf{Retrieval, ranking and LLMs}
\begin{checklist}
  \item Candidate generation and ranking stages
  \item Two-tower models and ANN search
  \item Cold start and position bias
  \item RAG pipeline and its failure points
  \item Evaluating LLM outputs
  \item Guardrails, cost and latency trade-offs
\end{checklist}
\end{minipage}

%----------------------------------------------------------------------
\needspace{16\baselineskip}
\section{Self-review rubric}\label{sec:rubric}

After each timed case, score yourself from 1 (missing) to 4 (strong) and write
one sentence on what to fix next time.

\begin{center}
\begin{tabularx}{\textwidth}{@{}l Y@{}}
\toprule
\textbf{Line} & \textbf{A score of 4 looks like} \\
\midrule
Problem framing & Asked for scale, latency and the business action before
  proposing anything \\
Metrics & Named offline, online and guardrail metrics and linked them \\
Data and labels & Raised label delay, noise and leakage unprompted \\
Modelling & Started simple, justified each step up in complexity \\
Serving & Gave a latency budget and a fallback path \\
Monitoring & Explained drift detection and the retraining trigger \\
Trade-offs & Stated at least three explicit trade-offs and chose a side \\
Communication & Kept to time, drew the system, checked in with the
  interviewer \\
\bottomrule
\end{tabularx}
\end{center}

%----------------------------------------------------------------------
\section{Common mistakes}

\begin{itemize}
  \item Jumping to the model before agreeing what the system is for.
  \item Spending 25 minutes on modelling and two on serving.
  \item Proposing a complex model with no baseline to compare against.
  \item Quoting AUC alone on an imbalanced problem.
  \item Describing a random train-test split for time-ordered data.
  \item Ignoring what happens when a feature source or the model is
    unavailable.
  \item Presenting one design as the only option instead of naming the
    trade-offs.
\end{itemize}

%----------------------------------------------------------------------
\section{If an interview arrives early}

A five-day compressed version, about 90 minutes per day:

\begin{enumerate}
  \item Memorise the framework; write out Case~1 in full.
  \item Huyen chapters 7 and 8; tick the Serving and Release checklists.
  \item Case~1 spoken and timed, then Case~2.
  \item One case outside fraud (Case~4) and, if the role mentions LLMs,
    Case~7.
  \item One full mock, scored against the rubric; fix the two lowest lines.
\end{enumerate}

%----------------------------------------------------------------------
\section{Resources}

\begin{itemize}
  \item Chip Huyen, \emph{Designing Machine Learning Systems} (O'Reilly, 2022).
    The main text for concepts.
  \item Ali Aminian and Alex Xu, \emph{Machine Learning System Design
    Interview} (2023). Worked cases in interview format.
  \item Stanford CS~329S, \emph{Machine Learning Systems Design}: lecture
    notes, free online, as a lighter alternative to the Huyen book.
  \item Engineering blogs from payments and marketplace companies on their
    fraud and risk platforms, for realistic architectures and vocabulary.
\end{itemize}

\end{document}